% chapter08.tex — Least Upper Bounds (source: ch08-least-upper-bounds.pdf).
% Book pp. 133–146. Slice pp. 2 and 4 are swapped (135 before 134); typed in
% book order. Appendix. Uniform Continuity is in this file after the problems.

\spivakchapter{8}{Least Upper Bounds}

This chapter reveals the most important property of the real numbers.
Nevertheless, it is merely a sequel to Chapter~7; the path which must be
followed has already been indicated, and further discussion would be useless
delay.

\begin{spivakdefn}
A set $A$ of real numbers is \textbf{bounded above} if there is a number $x$
such that
\[
x \geq a \qquad\text{for every $a$ in $A$.}
\]
Such a number $x$ is called an \textbf{upper bound} for $A$.
\end{spivakdefn}

Obviously $A$ is bounded above if and only if there is a number $x$ which is
an upper bound for $A$ (and in this case there will be lots of upper bounds
for $A$); we often say, as a concession to idiomatic English, that ``$A$ has
an upper bound'' when we mean that there is a number which is an upper bound
for $A$.

Notice that the term ``bounded above'' has now been used in two ways---first,
in Chapter~7, in reference to functions, and now in reference to sets. This
dual usage should cause no confusion, since it will always be clear whether
we are talking about a set of numbers or a function. Moreover, the two
definitions are closely connected: if $A$ is the set $\{f(x) : a \leq x \leq b\}$,
then the function $f$ is bounded above on $[a,b]$ if and only if the set $A$
is bounded above.

The entire collection $\mathbf{R}$ of real numbers, and the natural numbers
$\mathbf{N}$, are both examples of sets which are \emph{not} bounded above.
An example of a set which \emph{is} bounded above is
\[
A = \{x : 0 \leq x < 1\}.
\]
To show that $A$ is bounded above we need only name some upper bound for $A$,
which is easy enough; for example, $138$ is an upper bound for $A$, and so are
$2$, $1\frac{1}{2}$, $1\frac{1}{4}$, and $1$. Clearly, $1$ is the least upper
bound of $A$; although the phrase just introduced is self-explanatory, in
order to avoid any possible confusion (in particular, to ensure that we all
know what the superlative of ``less'' means), we define this explicitly.

\begin{spivakdefn}
A number $x$ is a \textbf{least upper bound} of $A$ if
\begin{align*}
&(1) && \text{$x$ is an upper bound of $A$,}\\
\text{and}\quad &(2) && \text{if $y$ is an upper bound of $A$, then $x \leq y$.}
\end{align*}
\end{spivakdefn}

The use of the indefinite article ``a'' in this definition was merely a
concession to temporary ignorance. Now that we have made a precise
definition, it is easily seen that if $x$ and $y$ are both least upper bounds
of $A$, then $x = y$. Indeed, in this case
\begin{align*}
x &\leq y, && \text{since $y$ is an upper bound, and $x$ is a least upper bound,}\\
\text{and}\quad
y &\leq x, && \text{since $x$ is an upper bound, and $y$ is a least upper bound;}
\end{align*}
it follows that $x = y$. For this reason we speak of \emph{the} least upper
bound of $A$. The term \textbf{supremum} of $A$ is synonymous and has one
advantage. It abbreviates quite nicely to
\[
\sup A \qquad\text{(pronounced ``soup $A$'')}
\]
and saves us from the abbreviation
\[
\operatorname{lub} A
\]
(which is nevertheless used by some authors).

There is a series of important definitions, analogous to those just given,
which can now be treated more briefly. A set $A$ of real numbers is
\textbf{bounded below} if there is a number $x$ such that
\[
x \leq a \qquad\text{for every $a$ in $A$.}
\]
Such a number $x$ is called a \textbf{lower bound} for $A$. A number $x$ is
the \textbf{greatest lower bound} of $A$ if
\begin{align*}
&(1) && \text{$x$ is a lower bound of $A$,}\\
\text{and}\quad &(2) && \text{if $y$ is a lower bound of $A$, then $x \geq y$.}
\end{align*}
The greatest lower bound of $A$ is also called the \textbf{infimum} of $A$,
abbreviated
\[
\inf A;
\]
some authors use the abbreviation
\[
\operatorname{glb} A.
\]

One detail has been omitted from our discussion so far---the question of
which sets have at least one, and hence exactly one, least upper bound or
greatest lower bound. We will consider only least upper bounds, since the
question for greatest lower bounds can then be answered easily (Problem~2).

If $A$ is not bounded above, then $A$ has no upper bound at all, so $A$
certainly cannot be expected to have a least upper bound. It is tempting to
say that $A$ does have a least upper bound if it has \emph{some} upper bound,
but, like the principle of mathematical induction, this assertion can fail to
be true in a rather special way. If $A = \emptyset$, then $A$ is bounded
above. Indeed, any number $x$ is an upper bound for $\emptyset$:
\[
x \geq y \quad\text{for every $y$ in $\emptyset$}
\]
simply because there is no $y$ in $\emptyset$. Since \emph{every} number is
an upper bound for $\emptyset$, there is surely no least upper bound for
$\emptyset$. With this trivial exception however, our assertion is
true---and very important, definitely important enough to warrant
consideration of details. We are finally ready to state the last property of
the real numbers which we need.

\begin{quote}
\prop{13}\ \ (The least upper bound property) If $A$ is a set of real
numbers, $A \neq \emptyset$, and $A$ is bounded above, then $A$ has a least
upper bound.
\end{quote}

Property P13 may strike you as anticlimactic, but that is actually one of its
virtues. To complete our list of basic properties for the real numbers we
require no particularly abstruse proposition, but only a property so simple
that we might feel foolish for having overlooked it. Of course, the least
upper bound property is not really so innocent as all that; after all, it
does \emph{not} hold for the rational numbers $\mathbf{Q}$. For example, if
$A$ is the set of all rational numbers $x$ satisfying $x^2 < 2$, then there
is no \emph{rational} number $y$ which is an upper bound for $A$ and which is
less than or equal to every other \emph{rational} number which is an upper
bound for $A$. It will become clear only gradually how significant P13 is,
but we are already in a position to demonstrate its power, by supplying the
proofs which were omitted in Chapter~7.

\medskip
\noindent\textbf{THEOREM 7-1}\enspace
\textit{If $f$ is continuous on $[a,b]$ and $f(a) < 0 < f(b)$, then there is
some number $x$ in $[a,b]$ such that $f(x) = 0$.}

\begin{proof}
Our proof is merely a rigorous version of the outline developed at the end of
Chapter~7---we will locate the smallest number $x$ in $[a,b]$ with $f(x) = 0$.

Define the set $A$, shown in Figure~1, as follows:
\[
A = \{x : a \leq x \leq b,\ \text{and $f$ is negative on the interval $[a,x]$}\}.
\]
\begin{center}
\figtex{assets/plots/ch08-fig01}\\[0.6em]
\figlabel{1}
\end{center}
Clearly $A \neq \emptyset$, since $a$ is in $A$; in fact, there is some
$\delta > 0$ such that $A$ contains all points $x$ satisfying
$a \leq x < a + \delta$; this follows from Problem~6-16, since $f$ is
continuous on $[a,b]$ and $f(a) < 0$. Similarly, $b$ is an upper bound for
$A$ and, in fact, there is a $\delta > 0$ such that all points $x$ satisfying
$b - \delta < x \leq b$ are upper bounds for $A$; this also follows from
Problem~6-16, since $f(b) > 0$.

From these remarks it follows that $A$ has a least upper bound $\alpha$ and
that $a < \alpha < b$. We now wish to show that $f(\alpha) = 0$, by
eliminating the possibilities $f(\alpha) < 0$ and $f(\alpha) > 0$.

Suppose first that $f(\alpha) < 0$. By Theorem~6-3, there is a $\delta > 0$
such that $f(x) < 0$ for $\alpha - \delta < x < \alpha + \delta$ (Figure~2).
\begin{center}
\figtex{assets/plots/ch08-fig02}\\[0.6em]
\figlabel{2}
\end{center}
Now there is some number $x_0$ in $A$ which satisfies
$\alpha - \delta < x_0 < \alpha$ (because otherwise $\alpha$ would not be the
\emph{least} upper bound of $A$). This means that $f$ is negative on the
whole interval $[a,x_0]$. But if $x_1$ is a number between $\alpha$ and
$\alpha + \delta$, then $f$ is also negative on the whole interval
$[x_0,x_1]$. Therefore $f$ is negative on the interval $[a,x_1]$, so $x_1$
is in $A$. But this contradicts the fact that $\alpha$ is an upper bound for
$A$; our original assumption that $f(\alpha) < 0$ must be false.

Suppose, on the other hand, that $f(\alpha) > 0$. Then there is a number
$\delta > 0$ such that $f(x) > 0$ for $\alpha - \delta < x < \alpha + \delta$
(Figure~3).
\begin{center}
\figtex{assets/plots/ch08-fig03}\\[0.6em]
\figlabel{3}
\end{center}
Once again we know that there is an $x_0$ in $A$ satisfying
$\alpha - \delta < x_0 < \alpha$; but this means that $f$ is negative on
$[a,x_0]$, which is impossible, since $f(x_0) > 0$. Thus the assumption
$f(\alpha) > 0$ also leads to a contradiction, leaving $f(\alpha) = 0$ as the
only possible alternative.
\end{proof}

The proofs of Theorems~2 and~3 of Chapter~7 require a simple preliminary
result, which will play much the same role as Theorem~6-3 played in the
previous proof.

\begin{theorem}
If $f$ is continuous at $a$, then there is a number $\delta > 0$ such that
$f$ is bounded above on the interval $(a - \delta, a + \delta)$ (see
Figure~4).
\end{theorem}

\begin{center}
\figtex{assets/plots/ch08-fig04}\\[0.6em]
\figlabel{4}
\end{center}

\begin{proof}
Since $\lim_{x \to a} f(x) = f(a)$, there is, for every $\varepsilon > 0$, a
$\delta > 0$ such that, for all $x$,
\[
\text{if }\abs{x - a} < \delta,\text{ then }\abs{f(x) - f(a)} < \varepsilon.
\]
It is only necessary to apply this statement to some particular $\varepsilon$
(any one will do), for example, $\varepsilon = 1$. We conclude that there is
a $\delta > 0$ such that, for all $x$,
\[
\text{if }\abs{x - a} < \delta,\text{ then }\abs{f(x) - f(a)} < 1.
\]
It follows, in particular, that if $\abs{x - a} < \delta$, then
$f(x) - f(a) < 1$. This completes the proof: on the interval
$(a - \delta, a + \delta)$ the function $f$ is bounded above by $f(a) + 1$.
\end{proof}

It should hardly be necessary to add that we can now also prove that $f$ is
bounded below on some interval $(a - \delta, a + \delta)$, and, finally, that
$f$ is bounded on some open interval containing $a$.

A more significant point is the observation that if
$\lim_{x \to a^{+}} f(x) = f(a)$, then there is a $\delta > 0$ such that $f$
is bounded on the set $\{x : a \leq x < a + \delta\}$, and a similar
observation holds if $\lim_{x \to b^{-}} f(x) = f(b)$. Having made these
observations (and assuming that you will supply the proofs), we tackle our
second major theorem.

\medskip
\noindent\textbf{THEOREM 7-2}\enspace
\textit{If $f$ is continuous on $[a,b]$, then $f$ is bounded above on
$[a,b]$.}

\begin{proof}
Let
\[
A = \{x : a \leq x \leq b\text{ and $f$ is bounded above on $[a,x]$}\}.
\]
Clearly $A \neq \emptyset$ (since $a$ is in $A$), and $A$ is bounded above
(by $b$), so $A$ has a least upper bound $\alpha$. Notice that we are here
applying the term ``bounded above'' both to the set $A$, which can be
visualized as lying on the horizontal axis, and to $f$, i.e., to the sets
$\{f(y) : a \leq y \leq x\}$, which can be visualized as lying on the
vertical axis (Figure~5).
\begin{center}
\figtex{assets/plots/ch08-fig05}\\[0.6em]
\figlabel{5}
\end{center}

Our first step is to prove that we actually have $\alpha = b$. Suppose,
instead, that $\alpha < b$. By Theorem~1 there is $\delta > 0$ such that $f$
is bounded on $(\alpha - \delta, \alpha + \delta)$. Since $\alpha$ is the
least upper bound of $A$ there is some $x_0$ in $A$ satisfying
$\alpha - \delta < x_0 < \alpha$. This means that $f$ is bounded on
$[a,x_0]$. But if $x_1$ is any number with $\alpha < x_1 < \alpha + \delta$,
then $f$ is also bounded on $[x_0,x_1]$. Therefore $f$ is bounded on
$[a,x_1]$, so $x_1$ is in $A$, contradicting the fact that $\alpha$ is an
upper bound for $A$. This contradiction shows that $\alpha = b$. One detail
should be mentioned: this demonstration implicitly assumed that $a < \alpha$
[so that $f$ would be defined on some interval
$(\alpha - \delta, \alpha + \delta)$]; the possibility $a = \alpha$ can be
ruled out similarly, using the existence of a $\delta > 0$ such that $f$ is
bounded on $\{x : a \leq x < a + \delta\}$.

The proof is not quite complete---we only know that $f$ is bounded on
$[a,x]$ for every $x < b$, not necessarily that $f$ is bounded on $[a,b]$.
However, only one small argument needs to be added.

There is a $\delta > 0$ such that $f$ is bounded on
$\{x : b - \delta < x \leq b\}$. There is $x_0$ in $A$ such that
$b - \delta < x_0 < b$. Thus $f$ is bounded on $[a,x_0]$ and also on
$[x_0,b]$, so $f$ is bounded on $[a,b]$.
\end{proof}

To prove the third important theorem we resort to a trick.

\medskip
\noindent\textbf{THEOREM 7-3}\enspace
\textit{If $f$ is continuous on $[a,b]$, then there is a number $y$ in
$[a,b]$ such that $f(y) \geq f(x)$ for all $x$ in $[a,b]$.}

\begin{proof}
We already know that $f$ is bounded on $[a,b]$, which means that the set
\[
\{f(x) : x\text{ in }[a,b]\}
\]
is bounded. This set is obviously not $\emptyset$, so it has a least upper
bound $\alpha$. Since $\alpha \geq f(x)$ for $x$ in $[a,b]$ it suffices to
show that $\alpha = f(y)$ for some $y$ in $[a,b]$.

Suppose instead that $\alpha \neq f(y)$ for all $y$ in $[a,b]$. Then the
function $g$ defined by
\[
g(x) = \frac{1}{\alpha - f(x)}, \qquad x\text{ in }[a,b]
\]
is continuous on $[a,b]$, since the denominator of the right side is never
$0$. On the other hand, $\alpha$ is the least upper bound of
$\{f(x) : x\text{ in }[a,b]\}$; this means that
\[
\text{for every $\varepsilon > 0$ there is $x$ in $[a,b]$ with
$\alpha - f(x) < \varepsilon$.}
\]
This, in turn, means that
\[
\text{for every $\varepsilon > 0$ there is $x$ in $[a,b]$ with
$g(x) > 1/\varepsilon$.}
\]
But \emph{this} means that $g$ is not bounded on $[a,b]$, contradicting the
previous theorem.
\end{proof}

At the beginning of this chapter the set of natural numbers $\mathbf{N}$ was
given as an example of an unbounded set. We are now going to \emph{prove}
that $\mathbf{N}$ is unbounded. After the difficult theorems proved in this
chapter you may be startled to find such an ``obvious'' theorem winding up
our proceedings. If so, you are, perhaps, allowing the geometrical picture
of $\mathbf{R}$ to influence you too strongly. ``Look,'' you may say, ``the
real numbers look like
\begin{center}
\figtex{assets/geometric/ch08-fig00}
\end{center}
so every number $x$ is between two integers $n$, $n+1$ (unless $x$ is itself
an integer).'' Basing the argument on a geometric picture is not a proof,
however, and even the geometric picture contains an assumption: that if you
place unit segments end-to-end you will eventually get a segment larger than
any given segment. This axiom, often omitted from a first introduction to
geometry, is usually attributed (not quite justly) to Archimedes, and the
corresponding property for numbers, that $\mathbf{N}$ is not bounded, is
called the \emph{Archimedean property} of the real numbers. This property is
\emph{not} a consequence of P1--P12 (see reference [14] of the Suggested
Reading), although it does hold for $\mathbf{Q}$, of course. Once we have
P13 however, there are no longer any problems.

\begin{theorem}
$\mathbf{N}$ is not bounded above.
\end{theorem}

\begin{proof}
Suppose $\mathbf{N}$ were bounded above. Since $\mathbf{N} \neq \emptyset$,
there would be a least upper bound $\alpha$ for $\mathbf{N}$. Then
\[
\alpha \geq n \qquad\text{for all $n$ in $\mathbf{N}$.}
\]
Consequently,
\[
\alpha \geq n + 1 \qquad\text{for all $n$ in $\mathbf{N}$,}
\]
since $n+1$ is in $\mathbf{N}$ if $n$ is in $\mathbf{N}$. But this means that
\[
\alpha - 1 \geq n \qquad\text{for all $n$ in $\mathbf{N}$,}
\]
and \emph{this} means that $\alpha - 1$ is also an upper bound for
$\mathbf{N}$, contradicting the fact that $\alpha$ is the least upper bound.
\end{proof}

There is a consequence of Theorem~2 (actually an equivalent formulation)
which we have very often assumed implicitly.

\begin{theorem}
For any $\varepsilon > 0$ there is a natural number $n$ with $1/n < \varepsilon$.
\end{theorem}

\begin{proof}
Suppose not; then $1/n \geq \varepsilon$ for all $n$ in $\mathbf{N}$. Thus
$n \leq 1/\varepsilon$ for all $n$ in $\mathbf{N}$. But this means that
$1/\varepsilon$ is an upper bound for $\mathbf{N}$, contradicting Theorem~2.
\end{proof}

A brief glance through Chapter~6 will show you that the result of Theorem~3
was used in the discussion of many examples. Of course, Theorem~3 was not
available at the time, but the examples were so important that in order to
give them some cheating was tolerated. As partial justification for this
dishonesty we can claim that this result was never used in the proof of a
\emph{theorem}, but if your faith has been shaken, a review of all the proofs
given so far is in order. Fortunately, such deception will not be necessary
again. We have now stated every property of the real numbers that we will
ever need. Henceforth, no more lies.

\subsection*{PROBLEMS}

\pnum{1.} Find the least upper bound and the greatest lower bound (if they
exist) of the following sets. Also decide which sets have greatest and least
elements (i.e., decide when the least upper bound and greatest lower bound
happens to belong to the set).
\begin{romans}
\item $\displaystyle\left\{\frac{1}{n} : n\text{ in }\mathbf{N}\right\}$.
\item $\displaystyle\left\{\frac{1}{n} : n\text{ in }\mathbf{Z}\text{ and }n \neq 0\right\}$.
\item $\{x : x = 0\text{ or }x = 1/n\text{ for some $n$ in }\mathbf{N}\}$.
\item $\{x : 0 \leq x \leq \sqrt{2}\text{ and $x$ is rational}\}$.
\item $\{x : x^2 + x + 1 \geq 0\}$.
\item $\{x : x^2 + x - 1 < 0\}$.
\item $\{x : x < 0\text{ and }x^2 + x - 1 < 0\}$.
\item $\displaystyle\left\{\frac{1}{n} + (-1)^n : n\text{ in }\mathbf{N}\right\}$.
\end{romans}

\pnum{2.}
\begin{letters}
\item Suppose $A \neq \emptyset$ is bounded below. Let $-A$ denote the set of
  all $-x$ for $x$ in $A$. Prove that $-A \neq \emptyset$, that $-A$ is
  bounded above, and that $-\sup(-A)$ is the greatest lower bound of $A$.
\item If $A \neq \emptyset$ is bounded below, let $B$ be the set of all lower
  bounds of $A$. Show that $B \neq \emptyset$, that $B$ is bounded above, and
  that $\sup B$ is the greatest lower bound of $A$.
\end{letters}

\pnum{3.} Let $f$ be a continuous function on $[a,b]$ with
$f(a) < 0 < f(b)$.
\begin{letters}
\item The proof of Theorem~7-1 showed that there is a smallest $x$ in
  $[a,b]$ with $f(x) = 0$. If there is more than one $x$ in $[a,b]$ with
  $f(x) = 0$, is there necessarily a second smallest? Show that there is a
  largest $x$ in $[a,b]$ with $f(x) = 0$. (Try to give an easy proof by
  considering a new function closely related to $f$.)
\item The proof of Theorem~7-1 depended upon considering
  $A = \{x : a \leq x \leq b$ and $f$ is negative on $[a,x]\}$. Give
  another proof of Theorem~7-1, which depends upon consideration of
  $B = \{x : a \leq x \leq b$ and $f(x) < 0\}$. Which point $x$ in
  $[a,b]$ with $f(x) = 0$ will this proof locate? Give an example where the
  sets $A$ and $B$ are not the same.
\end{letters}

\pnum{*4.}
\begin{letters}
\item Suppose that $f$ is continuous on $[a,b]$ and that $f(a) = f(b) = 0$.
  Suppose also that $f(x_0) > 0$ for some $x_0$ in $[a,b]$. Prove that there
  are numbers $c$ and $d$ with $a \leq c < x_0 < d \leq b$ such that
  $f(c) = f(d) = 0$, but $f(x) > 0$ for all $x$ in $(c,d)$. Hint: The
  previous problem can be used to good advantage.
\item Suppose that $f$ is continuous on $[a,b]$ and that $f(a) < f(b)$.
  Prove that there are numbers $c$ and $d$ with $a \leq c < d \leq b$ such
  that $f(c) = f(a)$ and $f(d) = f(b)$ and $f(a) < f(x) < f(d)$ for all $x$
  in $(c,d)$.
\end{letters}

\pnum{5.}
\begin{letters}
\item Suppose that $y - x > 1$. Prove that there is an integer $k$ such that
  $x < k < y$. Hint: Let $l$ be the largest integer satisfying $l \leq x$,
  and consider $l + 1$.
\item Suppose $x < y$. Prove that there is a rational number $r$ such that
  $x < r < y$. Hint: If $1/n < y - x$, then $ny - nx > 1$. (Query: Why have
  parts (a) and (b) been postponed until this problem set?)
\item Suppose that $r < s$ are rational numbers. Prove that there is an
  irrational number between $r$ and $s$. Hint: As a start, you know that
  there is an irrational number between $0$ and $1$.
\item Suppose that $x < y$. Prove that there is an irrational number between
  $x$ and $y$. Hint: It is unnecessary to do any more work; this follows from
  (b) and (c).
\end{letters}

\pnum{*6.} A set $A$ of real numbers is said to be \textbf{dense} if every
open interval contains a point of $A$. For example, Problem~5 shows that the
set of rational numbers and the set of irrational numbers are each dense.
\begin{letters}
\item Prove that if $f$ is continuous and $f(x) = 0$ for all numbers $x$ in a
  dense set $A$, then $f(x) = 0$ for all $x$.
\item Prove that if $f$ and $g$ are continuous and $f(x) = g(x)$ for all $x$
  in a dense set $A$, then $f(x) = g(x)$ for all $x$.
\item If we assume instead that $f(x) \geq g(x)$ for all $x$ in $A$, show
  that $f(x) \geq g(x)$ for all $x$. Can $\geq$ be replaced by $>$ throughout?
\end{letters}

\pnum{7.} Prove that if $f$ is continuous and $f(x+y) = f(x) + f(y)$ for all
$x$ and $y$, then there is a number $c$ such that $f(x) = cx$ for all $x$.
(This conclusion can be demonstrated simply by combining the results of two
previous problems.) Point of information: There \emph{do} exist
\emph{noncontinuous} functions $f$ satisfying $f(x+y) = f(x) + f(y)$ for all
$x$ and $y$, but we cannot prove this now; in fact, this simple question
involves ideas that are usually never mentioned in undergraduate courses
(see reference [7] in the Suggested Reading).

\pnum{*8.} Suppose that $f$ is a function such that $f(a) \leq f(b)$ whenever
$a < b$ (Figure~6).
\begin{center}
\figtex{assets/plots/ch08-fig06}\\[0.6em]
\figlabel{6}
\end{center}
\begin{letters}
\item Prove that $\lim_{x \to a^{-}} f(x)$ and $\lim_{x \to a^{+}} f(x)$ both
  exist. Hint: Why is this problem in this chapter?
\item Prove that $f$ never has a removable discontinuity (this terminology
  comes from Problem~6-17).
\item Prove that if $f$ satisfies the conclusions of the Intermediate Value
  Theorem, then $f$ is continuous.
\end{letters}

\pnum{*9.} If $f$ is a bounded function on $[0,1]$, let
$\lvert\lvert\lvert f \rvert\rvert\rvert
= \sup\{\abs{f(x)} : x\text{ in }[0,1]\}$.
Prove analogues of the properties of $\lVert\,\rVert$ in Problem~7-14.

\pnum{10.} Suppose $\alpha > 0$. Prove that every number $x$ can be written
uniquely in the form $x = k\alpha + x'$, where $k$ is an integer, and
$0 \leq x' < \alpha$.

\pnum{11.}
\begin{letters}
\item Suppose that $a_1, a_2, a_3, \dots$ is a sequence of positive numbers
  with $a_{n+1} \leq a_n/2$. Prove that for any $\varepsilon > 0$ there is
  some $n$ with $a_n < \varepsilon$.
\item Suppose $P$ is a regular polygon inscribed inside a circle. If $P'$ is
  the inscribed regular polygon with twice as many sides, show that the
  difference between the area of the circle and the area of $P'$ is less than
  half the difference between the area of the circle and the area of $P$
  (use Figure~7).
\begin{center}
\figtex{assets/geometric/ch08-fig07}\\[0.6em]
\figlabel{7}
\end{center}
\item Prove that there is a regular polygon $P$ inscribed in a circle with
  area as close as desired to the area of the circle. In order to do part~(c)
  you will need part~(a). This was clear to the Greeks, who used part~(a) as
  the basis for their entire treatment of proportion and area. By calculating
  the areas of polygons, this method (``the method of exhaustion'') allows
  computations of $\pi$ to any desired accuracy; Archimedes used it to show
  that $\frac{223}{71} < \pi < \frac{22}{7}$. But it has far greater
  theoretical importance:
\item[*(d)] Using the fact that the areas of two regular polygons with the
  same number of sides have the same ratio as the square of their sides,
  prove that the areas of two circles have the same ratios as the square of
  their radii. Hint: Deduce a contradiction from the assumption that the
  ratio of the areas is greater, or less, than the ratio of the square of the
  radii by inscribing appropriate polygons.
\end{letters}

\pnum{12.} Suppose that $A$ and $B$ are two nonempty sets of numbers such that
$x \leq y$ for all $x$ in $A$ and all $y$ in $B$.
\begin{letters}
\item Prove that $\sup A \leq y$ for all $y$ in $B$.
\item Prove that $\sup A \leq \inf B$.
\end{letters}

\pnum{13.} Let $A$ and $B$ be two nonempty sets of numbers which are bounded
above, and let $A + B$ denote the set of all numbers $x + y$ with $x$ in $A$
and $y$ in $B$. Prove that $\sup(A+B) = \sup A + \sup B$. Hint: The
inequality $\sup(A+B) \leq \sup A + \sup B$ is easy. Why? To prove that
$\sup A + \sup B \leq \sup(A+B)$ it suffices to prove that
$\sup A + \sup B \leq \sup(A+B) + \varepsilon$ for all $\varepsilon > 0$;
begin by choosing $x$ in $A$ and $y$ in $B$ with
$\sup A - x < \varepsilon/2$ and $\sup B - y < \varepsilon/2$.

\begin{center}
\figtex{assets/geometric/ch08-fig08}\\[0.6em]
\figlabel{8}
\end{center}
\pnum{14.}
\begin{letters}
\item Consider a sequence of closed intervals $I_1 = [a_1,b_1]$,
  $I_2 = [a_2,b_2]$,~\dots. Suppose that $a_n \leq a_{n+1}$ and
  $b_{n+1} \leq b_n$ for all $n$ (Figure~8). Prove that there is a point $x$
  which is in every $I_n$.
\item Show that this conclusion is false if we consider open intervals
  instead of closed intervals.
\end{letters}

The simple result of Problem~14(a) is called the ``Nested Interval Theorem.''
It may be used to give alternative proofs of Theorems~1 and~2. The
appropriate reasoning, outlined in the next two problems, illustrates a
general method, called a ``bisection argument.''

\pnum{*15.} Suppose $f$ is continuous on $[a,b]$ and $f(a) < 0 < f(b)$. Then
either $f((a+b)/2) = 0$, or $f$ has different signs at the end points of the
interval $[a,(a+b)/2]$, or $f$ has different signs at the end points of
$[(a+b)/2,b]$. Why? If $f((a+b)/2) \neq 0$, let $I_1$ be the interval on
which $f$ changes sign. Now bisect $I_1$. Either $f$ is $0$ at the midpoint,
or $f$ changes sign on one of the two intervals. Let $I_2$ be that interval.
Continue in this way, to define $I_n$ for each $n$ (unless $f$ is $0$ at some
midpoint). Use the Nested Interval Theorem to find a point $x$ where
$f(x) = 0$.

\pnum{*16.} Suppose $f$ were continuous on $[a,b]$, but not bounded on
$[a,b]$. Then $f$ would be unbounded on either $[a,(a+b)/2]$ or
$[(a+b)/2,b]$. Why? Let $I_1$ be one of these intervals on which $f$ is
unbounded. Proceed as in Problem~15 to obtain a contradiction.

\pnum{17.}
\begin{letters}
\item Let $A = \{x : x < \alpha\}$. Prove the following (they are all easy):
\begin{romans}
\item If $x$ is in $A$ and $y < x$, then $y$ is in $A$.
\item $A \neq \emptyset$.
\item $A \neq \mathbf{R}$.
\item If $x$ is in $A$, then there is some number $x'$ in $A$ such that
  $x < x'$.
\end{romans}
\item Suppose, conversely, that $A$ satisfies (i)--(iv). Prove that
  $A = \{x : x < \sup A\}$.
\end{letters}

\pnum{*18.} A number $x$ is called an \textbf{almost upper bound} for $A$ if
there are only finitely many numbers $y$ in $A$ with $y \geq x$. An
\textbf{almost lower bound} is defined similarly.
\begin{letters}
\item Find all almost upper bounds and almost lower bounds of the sets in
  Problem~1.
\item Suppose that $A$ is a bounded infinite set. Prove that the set $B$ of
  all almost upper bounds of $A$ is nonempty, and bounded below.
\item It follows from part~(b) that $\inf B$ exists; this number is called
  the \textbf{limit superior} of $A$, and denoted by $\varlimsup A$ or
  $\limsup A$. Find $\varlimsup A$ for each set $A$ in Problem~1.
\item Define $\varliminf A$, and find it for all $A$ in Problem~1.
\end{letters}

\pnum{*19.} If $A$ is a bounded infinite set prove
\begin{letters}
\item $\varliminf A \leq \varlimsup A$.
\item $\varlimsup A \leq \sup A$.
\item If $\varlimsup A < \sup A$, then $A$ contains a largest element.
\item The analogues of parts (b) and (c) for $\varliminf$.
\end{letters}

\begin{center}
\figtex{assets/plots/ch08-fig09}\\[0.6em]
\figlabel{9}
\end{center}
\pnum{20.} Let $f$ be a continuous function on $\mathbf{R}$. A point $x$ is
called a \textbf{shadow point} of $f$ if there is a number $y > x$ with
$f(y) > f(x)$. The rationale for this terminology is indicated in Figure~9;
the parallel lines are the rays of the sun rising in the east (you are facing
north). Suppose that all points of $(a,b)$ are shadow points, but that $a$
and $b$ are not shadow points. Clearly, $f(a) \geq f(b)$.
\begin{letters}
\item Suppose that $f(a) > f(b)$. Show that the point where $f$ takes on its
  maximum value on $[a,b]$ must be $a$.
\item Then show that this leads to a contradiction, so that in fact we must
  have $f(a) = f(b)$.
\end{letters}
This little result, known as the Rising Sun Lemma, is instrumental in proving
several beautiful theorems that do not appear in this book; see page~450.

\subsection*{APPENDIX.\quad UNIFORM CONTINUITY}

Now that we've come to the end of the ``foundations,'' it might be
appropriate to slip in one further fundamental concept. This notion is not
used crucially in the rest of the book, but it can help clarify many points
later on.

We know that the function $f(x) = x^2$ is continuous at $a$ for all $a$. In
other words,
\begin{quote}
if $a$ is any number, then for every $\varepsilon > 0$ there is some
$\delta > 0$ such that, for all $x$, if $\abs{x - a} < \delta$, then
$\abs{x^2 - a^2} < \varepsilon$.
\end{quote}
Of course, $\delta$ depends on $\varepsilon$. But $\delta$ also \emph{depends
on $a$}---the $\delta$ that works at $a$ might not work at $b$ (Figure~1).
\begin{center}
\figtex{assets/plots/ch08-fig10}\\[0.6em]
\figlabel{1}
\end{center}
Indeed, it's clear that given $\varepsilon > 0$ there is no one $\delta > 0$
that works for all $a$, or even for all positive $a$. In fact, the number
$a + \delta/2$ will certainly satisfy $\abs{x - a} < \delta$, but if
$a > 0$, then
\[
\abs[\bigg]{\biggl(a + \frac{\delta}{2}\biggr)^2 - a^2}
= \abs[\bigg]{a\delta + \frac{\delta^2}{4}} \geq a\delta,
\]
and this won't be $< \varepsilon$ once $a > \varepsilon/\delta$. (This is
just an admittedly confusing computational way of saying that $f$ is growing
faster and faster!)

On the other hand, for any $\varepsilon > 0$ there \emph{will} be one
$\delta > 0$ that works for all $a$ in any interval $[-N,N]$. In fact, the
$\delta$ which works at $N$ or $-N$ will also work everywhere else in the
interval.

As a final example, consider the function $f(x) = \sin 1/x$, or the function
whose graph appears in Figure~18 on page~62. It is easy to see that, so long
as $\varepsilon < 1$, there will not be one $\delta > 0$ that works for these
functions at all points $a$ in the open interval $(0,1)$.

These examples illustrate important distinctions between the behavior of
various continuous functions on certain intervals, and there is a special
term to signal this distinction.

\begin{spivakdefn}
The function $f$ is \textbf{uniformly continuous on an interval $A$} if for
every $\varepsilon > 0$ there is some $\delta > 0$ such that, for all $x$ and
$y$ in $A$,
\[
\text{if }\abs{x - y} < \delta,\text{ then }\abs{f(x) - f(y)} < \varepsilon.
\]
\end{spivakdefn}

We've seen that a function can be continuous on the whole line, or on an open
interval, without being uniformly continuous there. On the other hand, the
function $f(x) = x^2$ did turn out to be uniformly continuous on any closed
interval. This shouldn't be too surprising---it's the same sort of thing that
occurs when we ask whether a function is bounded on an interval---and we
would be led to suspect that any continuous function on a closed interval is
also uniformly continuous on that interval. In order to prove this, we'll
need to deal first with one subtle point.

Suppose that we have two intervals $[a,b]$ and $[b,c]$ with the common
end-point $b$, and a function $f$ that is continuous on $[a,c]$. Let
$\varepsilon > 0$ and suppose that the following two statements hold:
\begin{romans}
\item if $x$ and $y$ are in $[a,b]$ and $\abs{x - y} < \delta_1$, then
  $\abs{f(x) - f(y)} < \varepsilon$,
\item if $x$ and $y$ are in $[b,c]$ and $\abs{x - y} < \delta_2$, then
  $\abs{f(x) - f(y)} < \varepsilon$.
\end{romans}
We'd like to know if there is some $\delta > 0$ such that
$\abs{f(x) - f(y)} < \varepsilon$ whenever $x$ and $y$ are points in $[a,c]$
with $\abs{x - y} < \delta$. Our first inclination might be to choose
$\delta$ as the minimum of $\delta_1$ and $\delta_2$. But it is easy to see
what goes wrong (Figure~2): we might have $x$ in $[a,b]$ and $y$ in $[b,c]$,
and then neither (i) nor (ii) tells us anything about $\abs{f(x) - f(y)}$.
So we have to be a little more cagey, and also use continuity of $f$ at $b$.
\begin{center}
\figtex{assets/geometric/ch08-fig11}\\[0.6em]
\figlabel{2}
\end{center}

\medskip
\noindent\textbf{LEMMA}\enspace
\textit{Let $a < b < c$ and let $f$ be continuous on the interval $[a,c]$.
Let $\varepsilon > 0$, and suppose that statements \textup{(i)} and
\textup{(ii)} hold. Then there is a $\delta > 0$ such that, if $x$ and $y$
are in $[a,c]$ and $\abs{x - y} < \delta$, then
$\abs{f(x) - f(y)} < \varepsilon$.}

\begin{proof}
Since $f$ is continuous at $b$, there is a $\delta_3 > 0$ such that,
\[
\text{if }\abs{x - b} < \delta_3,\text{ then }
\abs{f(x) - f(b)} < \frac{\varepsilon}{2}.
\]
It follows that
\begin{quote}
(iii)\quad if $\abs{x - b} < \delta_3$ and $\abs{y - b} < \delta_3$, then
$\abs{f(x) - f(y)} < \varepsilon$.
\end{quote}
Choose $\delta$ to be the minimum of $\delta_1$, $\delta_2$, and $\delta_3$.
We claim that this $\delta$ works. In fact, suppose that $x$ and $y$ are any
two points in $[a,c]$ with $\abs{x - y} < \delta$. If $x$ and $y$ are both in
$[a,b]$, then $\abs{f(x) - f(y)} < \varepsilon$ by (i); and if $x$ and $y$
are both in $[b,c]$, then $\abs{f(x) - f(y)} < \varepsilon$ by (ii). The only
other possibility is that
\[
x < b < y \qquad\text{or}\qquad y < b < x.
\]
In either case, since $\abs{x - y} < \delta$, we also have
$\abs{x - b} < \delta$ and $\abs{y - b} < \delta$. So
$\abs{f(x) - f(y)} < \varepsilon$ by (iii).
\end{proof}

\setcounter{theorem}{0}
\begin{theorem}
If $f$ is continuous on $[a,b]$, then $f$ is uniformly continuous on
$[a,b]$.
\end{theorem}

\begin{proof}
It's the usual trick, but we've got to be a little bit careful about the
mechanism of the proof. For $\varepsilon > 0$ let's say that $f$ is
\emph{$\varepsilon$-good} on $[a,b]$ if there is some $\delta > 0$ such that,
for all $y$ and $z$ in $[a,b]$,
\[
\text{if }\abs{y - z} < \delta,\text{ then }\abs{f(y) - f(z)} < \varepsilon.
\]
Then we're trying to prove that $f$ is $\varepsilon$-good on $[a,b]$ for all
$\varepsilon > 0$.

Consider any particular $\varepsilon > 0$. Let
\[
A = \{x : a \leq x \leq b\text{ and $f$ is $\varepsilon$-good on $[a,x]$}\}.
\]
Then $A \neq \emptyset$ (since $a$ is in $A$), and $A$ is bounded above (by
$b$), so $A$ has a least upper bound $\alpha$. We really should write
$\alpha_{\varepsilon}$, since $A$ and $\alpha$ might depend on $\varepsilon$.
But we won't since we intend to prove that $\alpha = b$, no matter what
$\varepsilon$ is.

Suppose that we had $\alpha < b$. Since $f$ is continuous at $\alpha$, there
is some $\delta_0 > 0$ such that, if $\abs{y - \alpha} < \delta_0$, then
$\abs{f(y) - f(\alpha)} < \varepsilon/2$. Consequently, if
$\abs{y - \alpha} < \delta_0$ and $\abs{z - \alpha} < \delta_0$, then
$\abs{f(y) - f(z)} < \varepsilon$. So $f$ is surely $\varepsilon$-good on the
interval $[\alpha - \delta_0, \alpha + \delta_0]$. On the other hand, since
$\alpha$ is the least upper bound of $A$, it is also clear that $f$ is
$\varepsilon$-good on $[a, \alpha - \delta_0]$. Then the Lemma implies that
$f$ is $\varepsilon$-good on $[a, \alpha + \delta_0]$, so $\alpha + \delta_0$
is in $A$, contradicting the fact that $\alpha$ is an upper bound.

To complete the proof we just have to show that $\alpha = b$ is actually in
$A$. The argument for this is practically the same: Since $f$ is continuous
at $b$, there is some $\delta_0 > 0$ such that, if $b - \delta_0 < y < b$,
then $\abs{f(y) - f(b)} < \varepsilon/2$. So $f$ is $\varepsilon$-good on
$[b - \delta_0, b]$. But $f$ is also $\varepsilon$-good on
$[a, b - \delta_0]$, so the Lemma implies that $f$ is $\varepsilon$-good on
$[a,b]$.
\end{proof}

\subsection*{PROBLEMS}

\pnum{1.}
\begin{letters}
\item For which of the following values of $\alpha$ is the function
  $f(x) = x^{\alpha}$ uniformly continuous on $[0,\infty)$:
  $\alpha = 1/3$, $1/2$, $2$, $3$?
\item Find a function $f$ that is continuous and bounded on $(0,1]$, but not
  uniformly continuous on $(0,1]$.
\item Find a function $f$ that is continuous and bounded on $[0,\infty)$ but
  which is not uniformly continuous on $[0,\infty)$.
\end{letters}

\pnum{2.}
\begin{letters}
\item Prove that if $f$ and $g$ are uniformly continuous on $A$, then so is
  $f + g$.
\item Prove that if $f$ and $g$ are uniformly continuous and bounded on $A$,
  then $fg$ is uniformly continuous on $A$.
\item Show that this conclusion does not hold if one of them isn't bounded.
\item Suppose that $f$ is uniformly continuous on $A$, that $g$ is uniformly
  continuous on $B$, and that $f(x)$ is in $B$ for all $x$ in $A$. Prove that
  $g \circ f$ is uniformly continuous on $A$.
\end{letters}

\pnum{3.} Use a ``bisection argument'' (page~142) to give another proof of
Theorem~1.

\pnum{4.} Derive Theorem~7-2 as a consequence of Theorem~1.
